\subsection{Images of integral seminorm maps}
\label{subsec:integral-spectrum-images}

\begin{lemma}[Nonzero completed tensor fibres]
\label{spectrum:integral-nontrivial-fibre}
\lean{EnvelopeProperty.Berkovich.Spectrum.nontrivial_fibre}
\leanok
\leanok
\uses{spectrum:residue-field,spectrum:completed-residue-field}
For an integral homomorphism $A\to B$ of arbitrary commutative rings
and a seminorm point $x$ of $A$, suppose its support contains the kernel
of the homomorphism. Then $H(x)\otimes_A B$ is a nonzero ring.
\end{lemma}
\begin{proof}
The support is prime. Lying over gives a prime of $B$ contracting to
that support, so the ordinary residue-field fibre $\kappa(x)\otimes_A B$
is nonzero. Extension of scalars from the field $\kappa(x)$ to its
genuine completed field $H(x)$ is faithfully flat and preserves
nonzero algebras. Tensor cancellation identifies this nonzero extension
with $H(x)\otimes_A B$, giving the assertion.
\end{proof}

\begin{theorem}[Exact image of an integral pullback]
\label{spectrum:integral-image}
\lean{EnvelopeProperty.Berkovich.Spectrum.exists_comap_eq,
EnvelopeProperty.Berkovich.Spectrum.exists_comap_eq_of_isIntegral,
EnvelopeProperty.Berkovich.Spectrum.exists_comap_eq_iff_of_isIntegral,
EnvelopeProperty.Berkovich.Spectrum.range_comap_of_isIntegral,
EnvelopeProperty.Berkovich.Spectrum.comap_surjective_of_isIntegral}
\leanok
\leanok
\uses{spectrum:integral-nontrivial-fibre,spectrum:integral-completed-fibres,
spectrum:quotient-zero-locus}
For an integral ring homomorphism $\varphi:A\to B$, a point $x\in\mathcal S(A)$
is the restriction of a point of $\mathcal S(B)$ if and only if
$\ker\varphi\subseteq\operatorname{supp}(x)$. Thus the image of
the integral pullback is the analytic zero locus of $\ker\varphi$.
If $\varphi$ is also injective, its seminorm pullback is surjective.
\end{theorem}
\begin{proof}
A restricted point vanishes on the kernel by direct evaluation.
Conversely, the preceding lemma makes the completed tensor fibre
nonzero, hence it has a maximal ideal. The genuine completed-fibre
classification supplies a seminorm point lying over $x$ with that
maximal tensor support. This proves existence with the exact obstruction.
The pointwise equivalence proves the image formula. For injective
$\varphi$ the kernel is zero, contained in every support, proving
surjectivity. This uses actual integral norm-extension existence;
no unrestricted scalar base-change existence is inferred.
\end{proof}

\begin{corollary}[Images with an arbitrary prescribed base]
\label{spectrum:relative-integral-image}
\lean{EnvelopeProperty.Berkovich.RelativeSpectrum.exists_comap_eq_of_isIntegral,
EnvelopeProperty.Berkovich.RelativeSpectrum.exists_comap_eq_iff_of_isIntegral,
EnvelopeProperty.Berkovich.RelativeSpectrum.range_comap_of_isIntegral,
EnvelopeProperty.Berkovich.RelativeSpectrum.comap_surjective_of_isIntegral}
\leanok
\leanok
\uses{spectrum:integral-image,spectrum:relative}
For an integral algebra homomorphism over any commutative ring carrying
any prescribed real seminorm, the image consists exactly of the relative
points whose support contains its kernel. In particular an injective
integral algebra homomorphism induces a surjective relative pullback.
\end{corollary}
\begin{proof}
Use the absolute existence theorem over the underlying point $x$.
The resulting $y$ has the prescribed base values automatically:
on every base coefficient, algebra-map compatibility and $y|_A=x$
identify its value with the prescribed value of $x$. Bundle these
equalities to obtain the relative point. Necessity is unchanged
under forgetting the base. The image formula and injective case
then follow exactly as in the absolute construction.
\end{proof}
